\subsection{Exact minimum-cost realization of the static axis--scale operator}

The reduction in trainable parameters does not by itself guarantee low execution
latency, because the branch-form implementation evaluates one depthwise axial
convolution for every enabled axis--scale pair. We therefore derive an exact
minimum-cost realization of the same learned DT1D operator, without changing its
parameters, training objective, routing rule, or hypothesis class.

For axis $u$, dilation $d_s$, and normalized shifted kernel
$\widetilde{k}_{u,s}$, define
\begin{equation}
(\mathcal{T}_{u,d_s,\widetilde{k}_{u,s}}x)[n]
=
\sum_{r=-R}^{R}\widetilde{k}_{u,s}[r]
 x[n+d_s r e_u],
\qquad R=M+1.
\end{equation}
The original static-router realization is
\begin{equation}
\mathcal{F}_{\mathrm{branch}}x
=
\sum_{u\in\mathcal{A}}\sum_{s=1}^{S}
\pi_{u,s}\mathcal{T}_{u,d_s,\widetilde{k}_{u,s}}x.
\end{equation}
For a boundary-compatible scale group $G$, define the aggregated Laurent kernel
\begin{equation}
K_{u,G}[t]
=
\sum_{s\in G}\sum_{r=-R}^{R}
\pi_{u,s}\widetilde{k}_{u,s}[r]
\mathbf{1}_{\{t=d_s r\}}.
\end{equation}
Then
\begin{equation}
\mathcal{T}_{u,1,K_{u,G}}x
=
\sum_{s\in G}
\pi_{u,s}\mathcal{T}_{u,d_s,\widetilde{k}_{u,s}}x,
\end{equation}
so the aggregation is exact rather than approximate.

A singleton group retains the original compact dilated kernel and has cost
$c(\{s\})=2M+3$. A multi-scale group has dense unit-dilation support cost
\begin{equation}
c(G)=2(M+1)\max_{s\in G}d_s+1,
\qquad |G|>1.
\end{equation}
We select the lexicographically optimal partition
\begin{equation}
\mathcal{P}^{\star}
=
\arg\min_{\mathcal{P}}^{\mathrm{lex}}
\left(
\sum_{G\in\mathcal{P}}c(G),
|\mathcal{P}|
\right),
\end{equation}
where the first component is the dense axial tap count and the second is the
number of axial convolution calls. After sorting scales by dilation, the optimum
is obtained by the dynamic-programming recurrence
\begin{equation}
J(i)=\min_{0\le j<i}
\left[J(j)+c(\{j,\ldots,i-1\})\right].
\end{equation}
Because the all-singleton partition is feasible, the selected realization cannot
increase the modeled dense tap count, and its number of groups cannot exceed the
original number of scales.

For the default $M=1$ and $\mathcal{D}=\{1,2,4\}$, the original per-axis cost is
$5+5+5=15$ taps and three calls. Complete fusion costs 17 taps. The optimum is
$\{\{1,2\},\{4\}\}$, with $9+5=14$ taps and two calls. Across height and width,
this changes the realization from six to four axial calls and from 30 to 28 dense
taps, while preserving the learned mapping.

The existing stability result is unchanged. In fact,
\begin{equation}
\sum_{u,G}\lVert K_{u,G}\rVert_1
\le
\sum_{u,s}\pi_{u,s}\lVert\widetilde{k}_{u,s}\rVert_1
\le 1,
\end{equation}
and hence $\lVert\mathcal{F}x\rVert_p\le\lVert x\rVert_p$. The spectral response
is also preserved because
\begin{equation}
\widehat{K}_{u,G}(\omega)
=
\sum_{s\in G}\pi_{u,s}
\widehat{\widetilde{k}}_{u,s}(d_s\omega).
\end{equation}
Thus, the cosine-modulated multiplier interpretation and the Fredholm-type
motivation remain valid.

We verified the realization against the uploaded branch implementation over 720
configurations covering multiple radii, dilation sets, padding modes, feature-map
sizes, floating-point types, and pointwise-block settings. The maximum output
and parameter-gradient differences were $4.77\times10^{-7}$ and
$1.19\times10^{-7}$, respectively, with 100\% prediction agreement and no cost
regression. A 50-step AdamW trajectory had zero maximum loss difference and a
maximum final parameter difference of $1.19\times10^{-7}$. Controlled CPU tests,
with caching disabled and only the realization switch changed, showed adapter-level
inference speedups from $1.14\times$ to $1.25\times$ and forward--backward speedups
from $1.23\times$ to $1.48\times$. These measurements demonstrate an exact
mathematical execution optimization rather than a change to the data, optimizer,
baselines, or evaluation pipeline.
